% ===== §11 The Beauty of Practice =====
\section{The Beauty of Practice}
\label{p22:sec:beauty-practice}

\noindent
Perception discerns the field's direction; practice sends the path that way. The beauty of practice is the output face of the sensibility, the generation of appropriateness in the concrete act, the gesture or placement or word that, in this state of the field, makes the coupling gain rather than spin. If perception is the registration of where the field stands, practice is the intervention that alters where it stands, and the two are distinct enough that a person may have one without the other, may perceive the field's flattening with great fineness and be unable to act so as to arrest it, or may act effectively without perceiving finely what the action does. This section develops the beauty of practice and draws, as the method requires, the analysis of the contradiction between perception and practice, which is the contradiction this face most directly carries, and which has a definite dominant aspect in the intimate field.

\subsection{Practice: the generation of appropriateness}
\label{p22:sub:practice-generation}

To practise beautifully is to generate appropriateness in the act, and this is a different accomplishment from perceiving it. Appropriateness, the section on the coupling mechanism established, is the signature of a field accumulating in the generative direction, and to generate it is to perform the act that, in this state, sends the field that way. The accomplishment is concrete and particular, since the appropriate act is a property of the token field in its token state, and so practice cannot be the application of a known technique any more than perception could be the application of a rule; it is the production, in the concrete, of the gesture this field now requires, which must be found freshly because the state is fresh. This is why practice is a beauty and not a skill in the engineering sense, why it is cultivated by exercise rather than learned by instruction: what is raised is not a repertoire of techniques but a responsiveness, the capacity to produce, in the concrete and unrepeatable case, the act that builds this field now. The cultivation of practice is the cultivation of that responsiveness, the training of the hand and the word to find the appropriate act as the field presents itself, and it is beautiful in the same broad sense the whole paper has used, the sense in which the production of the fitting where no rule specifies it is the work of an aesthetic rather than a technical power.

The point bears comparison with the acquisition of any embodied skill that meets unrepeatable situations, the improviser's, the physician's, the craftsman's, because the comparison shows both what practice shares with skill and where it exceeds the engineering sense of the word. The expert improviser does not apply a rule for the next phrase; she has, through long exercise, acquired a responsiveness that produces the apt phrase as the music presents its state, and if asked for the rule she followed she could give none, because there was none, only a trained capacity to meet the concrete moment with the fitting act. This is practice in the sense the section means, and the analogy holds far. But it must be qualified in one respect that marks where the beauty of practice exceeds skill: the improviser's responsiveness is to a musical state, a thing without a face, whereas the field-builder's responsiveness is to a relational state that includes another generative subject, and the apt act is apt not to a passive material but to a person whose own perception and reception answer it. The training of practice is therefore the training of a responsiveness to a state that answers back, a skill whose material is a second freedom, and this is why its cultivation cannot be the cultivation of a craft upon an inert medium but must be the cultivation of a responsiveness within a coupling, learned only in the exercise of coupling and not in any practice upon a material that does not answer. The improviser trains alone with the instrument; the field-builder can train only with the other, because the state her practice must meet is a state the other co-constitutes, which returns the cultivation of practice, like everything in this paper, to the irreducibly relational.

This has a consequence for how practice fails, distinct from how perception fails, and worth marking because the dialectical analysis will turn on it. Perception fails by not seeing, the field's state going unregistered; practice fails by not reaching, the field's state seen but the apt act not produced, the responsiveness lacking or misfiring so that what is perceived to be needed is not delivered. The two failures are independent, and their independence is the content of the perception-practice contradiction the next part analyses. A partner may see with perfect clarity that the other needs a particular meeting and yet be unable to produce it, the hand or the word failing where the perception succeeded; another may produce apt acts readily while perceiving only coarsely what they are for. The cultivation of practice addresses the second failure specifically, the failure to reach, and it is a cultivation distinct from the cultivation of perception, the training of a responsiveness rather than a resolution, the raising of the capacity to deliver the apt act and not only to discern it.

\subsection{Practice is the dominant aspect, and its transformation}
\label{p22:sub:practice-dominant}

The contradiction between perception and practice is the secondary contradiction the method named first, and within it, in the intimate field, practice is the dominant aspect. The reason is that a relation is built by what is done and not by what is seen; the holonomy of a field accumulates through the acts that send its path the generative way, and a perception however fine that issues in no act adds nothing to the field's accumulation. A breakfast made with full attention builds the field; a fine perception of the other's state that produces no answering act leaves the field where it was. Perception matters, but it matters as it issues in practice, its value finally cashed in the act it informs, and so in the contradiction between them practice is dominant, the side on which the field's accumulation actually turns. This is a correction the tradition needs, since the tradition's emphasis on the cultivation of perceptive taste can leave the impression that to perceive finely is the whole of the aesthetic accomplishment, whereas in the building of a relational field the perceiving is subordinate to the doing, the discernment a preparation for the act in which alone the field is built.

\claim{In the contradiction between perception and practice, practice is the dominant aspect in the intimate field, because a relation accumulates through what is done and not through what is perceived. A fine perception that issues in no act leaves the field where it was; the holonomy accumulates in the gesture that sends the path the generative way. Perception is subordinate in the precise sense that its value is cashed in the practice it informs. This dominance transforms, however, when much is done and done wrongly: in a field where action is copious but blind, where the partners act constantly and constantly mis-act for want of perceiving what their acts do, the dominant aspect shifts to perception, the recovery of the seeing that would let the abundant action find its mark.}

The transformation is as important as the dominance, and it is the instance the method anticipated when it said the dominant and subordinate aspects reverse as the field's state changes. Ordinarily, in the intimate field, practice dominates, and the cultivation should weight the doing, the production of the appropriate act, over the refinement of a perception that risks becoming a connoisseurship of the relation's states without intervention in them. But when a field is full of action that misses, when the partners do much and build little because they act without perceiving what their action does to the field, the dominance reverses, and the work shifts to perception, since here the limiting factor is not the will to act but the sight to act well. A field of blind copious action is not helped by more action; it is helped by the perception that would let the action find its mark, and the dialectical reading is what tells the field-builder which case obtains, whether the field wants more doing or better seeing, a question answered not by a rule but by reading which aspect of the perception-practice contradiction now dominates the field's failure to accumulate.

\subsection{Against the over-correction toward doing}
\label{p22:sub:practice-overcorrection}

The claim that practice dominates carries a risk of over-correction that the section must guard against, lest the correction of the tradition's over-valuation of perception install the opposite error, an over-valuation of doing that mistakes mere activity for field-building. The dominance of practice is the dominance of the apt act, the gesture that, in this state, sends the path the generative way, and not the dominance of activity as such; a field is built by the right doing and not by much doing, and a partner who responds to the claim of practice's dominance by simply doing more, by multiplying gestures without the perception that would make them apt, has not built the field but only crowded it. The dominance of practice presupposes the perception that practice is subordinate to in the order of dependence, since the apt act is apt to a perceived state, and so the dominance of practice is not the license to dispense with perception but the insistence that perception be cashed in apt action rather than hoarded as connoisseurship. This is the fine point the dialectical reading must hold: practice dominates because the field is built by doing, and yet the doing that builds is the apt doing that perception informs, so that to weight practice is not to abandon perception but to insist that perception issue in the act where the field is actually built. The over-correction toward mere doing is as much a failure as the tradition's over-valuation of mere perceiving, and the beauty of practice is neither, being the production of the apt act in which a perceived state is met by a gesture that builds, the doing informed by seeing and cashed in the field's accumulation.

This guards, too, against a misreading of the cultivation that the worked example might otherwise invite, the misreading on which to cultivate practice is to drill a repertoire of helpful gestures, to learn the things one does to build a field as one might learn a set of techniques. The beauty of practice is not a repertoire, because the apt act is a property of the token state and no repertoire of types reaches it, and a partner armed with a repertoire of field-building gestures would apply them to states they do not fit, doing the helpful thing at the unhelpful moment, the repertoire substituting for the responsiveness that alone produces the apt act. The cultivation of practice is the cultivation of responsiveness and not the acquisition of a repertoire, the training of the capacity to produce the act this state requires and not the stocking of a kit of acts that worked before, and the dominance of practice is the dominance of this responsiveness, the trained capacity to do the apt thing, over a perception that would discern without doing, and not the dominance of a repertoire over a responsiveness or of activity over perception.
